% Nouveau EGA V (cote 25): archive cote 25, site page 74 (PDF page 75).
% typescript. Draft compared: none.
% AI-transcribed by Claude Opus 5.5 (claude-opus-5-5, adaptive thinking, effort high), Message Batches API, batch ega25b (msgbatch_01A7n5YuXubUPNK1zsepi8RG), 2026-11-13; not yet checked by eye.
\scanpage{25/74}
\mspage{\hand{\unsure{61|6}}}

sous-préschéma fermé d e Grass$^{m}$ , en particulier on peut former $\underline{X}^{(m)}_{Q}$ ,
qu'on se propose d'étudier. Un premier point, qui doit en tout état de
cause figurer dans le texte, c'est que $\underline{X}^{(m)}_{Q} \longrightarrow X$ est encore birationnel
du moins lorsque C coupe ``régulièrement'' X , et de façon précise $\underline{X}^{(m)}_{Q}$
est dans œ cas canoniquement isomorphe \unsure{àu} préschéma déduit de X en
faisant éclater $X\mathrm{x}_{P}\unsure{C}$ : la démonstration de ce fait n'est autre que celle
de 10.2. via \over{10}{\unsure{18}}.3. Un deuxième point,qui présente un certain intérêt,
mais que nous ne sommes pas absolument obligés d'inclure\xout{r} consiste à
dire que si on a choisi C ``assez général'', alors $\underline{X}^{(m)}_{Q} \longrightarrow Q$ a certaines
propriétés sympathiques, la plus classique étant celle-ci: X étant
supposé lisse sur S=Spec(k)\inserted[typed]{de dim $n \geq m$},et \xout{toutes}\inserted[typed]{propre et} géom irréd , Alors pour C ``assez
général'', l'ensemble\inserted[typed]{T} des $\hand{\zeta}\in Q$ pour lesquels $\underline{X}^{(m)}$ n'est\inserted[typed]{pas} lisse de dimen-%
sion n-m sur k($\hand{\zeta}$) est géom irréductible sur k($\hand{\zeta}$)\inserted[typed]{et de codim 1 dans Q}, et \xout{pour}
l'ensemble\inserted[typed]{T$^{1}$} des $\hand{\zeta}\in T$ pour lesquels $\underline{X}^{(m)}$ est ``supersingulier'' en au
moins un point est \xout{r}rare dans T;\xout{\unsure{ainsi, pour tout $\zeta \in T-T^{1}$, $\underline{X}^{(m)}_{\zeta}$ est}}
\xout{\unsure{n'admet que des points de}} enfin , si f:X$\rightarrow$P est une immersion, alors
quitte à agrandir un peu T$^{1}$ , pour tout $\hand{\zeta}\in$ T-T$^{1}$ , il y a exactement un
\margin{$|$}point non lisse dans $X^{(m)}_{\hand{\zeta}}$ , et ce dernier est rationnel sur k($\hand{\zeta}$). J'ai
oublié dans l'énoncé de spécifier qu'on supposait X \xout{de} $\longrightarrow$ P non ramifié,
et que l'on devait au préalable remplacer f par $\phi_{n}f$ , n $\geq$2 (où $\phi_{n}$ est
défini dans 9.1.). Le plus naturel \xout{semble} pour prouver cet énoncé sem-%
ble d'utiliser le sous-schéma Z (noté T dans 8.8.) des Grass$^{m}$ tels
que $\underline{X}^{(m)}$ soit \xout{``supers}Singulier\unsure{''}: on sait, sous les conditions dites, que Z
est géom. irréd. de codimension 1, et que le sous-schéma Z$^{1}$ correspondant
à des $\underline{X}^{(m)}$ ``supersinguliers'' y est r are. Il reste\inserted{rait} alors à prouver un
lemme de la nature suivante: soit Z une partie fermée de Grass$^{m}$ de codi-%
mension q, alors définissant Q(C) en t ermes de C comme ci-dessus, pour
tout C ``assez général'' l'intersection $Q(\unsure{Z|C})\cap Z$ est de codim $\geq$ q dans
\inserted{$Q(C)$; de plus, si $Z$ géom. irréd., itou $Q(C)\cap Z$ si $Z$ ``assez général''.}

\mspage{\hand{74}}
